
%====================== 字体设置（中英文） ======================%
% 中文字体：
%   - main.tex 使用 fontset=fandol（TeX Live 自带），跨平台最稳定。
%   - 因此这里不再手动 \setCJKmainfont，避免与 ctex fontset 冲突。
%
% 英文字体（建议使用 TeX Live 自带字体，避免系统缺字导致编译失败）：
%   - Times 风格：TeX Gyre Termes
%   - Helvetica 风格：TeX Gyre Heros
%   - Courier 风格：TeX Gyre Cursor
%
% 若你在 Windows 且希望使用 Times New Roman/Arial/Courier New，可将下面三行替换为：
  % \setmainfont{Times New Roman}
  % \setsansfont{Arial}
  % \setmonofont{Courier New}

% \setmainfont{TeX Gyre Termes}
% \setsansfont{TeX Gyre Heros}
% \setmonofont{TeX Gyre Cursor}

%====================== 英文字体：自动适配 ======================%
% \usepackage{fontspec}

% \IfFontExistsTF{Times New Roman}{
%   \setmainfont{Times New Roman}
% }{
%   \setmainfont{TeX Gyre Termes}
% }

% \IfFontExistsTF{Arial}{
%   \setsansfont{Arial}
% }{
%   \setsansfont{TeX Gyre Heros}
% }

% \IfFontExistsTF{Courier New}{
%   \setmonofont{Courier New}
% }{
%   \setmonofont{TeX Gyre Cursor}
% }

% 加载文件包中的字体
% =========================
% Local fonts config (./fonts)
% Compile with XeLaTeX or LuaLaTeX
% =========================
\usepackage{fontspec}



% 让 fontspec 可以直接按文件名找字体（不依赖系统安装）
\defaultfontfeatures{Ligatures=TeX}

% ---------- 英文字体（来自 ./fonts） ----------
\setmainfont{TimesNewRoman.ttf}[
  Path = ./fonts/,
  UprightFont = TimesNewRoman.ttf,
  BoldFont    = timesbd.ttf,
  ItalicFont  = timesi.ttf,
  BoldItalicFont = timesbi.ttf
]
% 如果你后续需要英文无衬线，可再补 Arial.ttf；目前你目录里没有，就先不设
% \setsansfont{Arial.ttf}[Path=./fonts/]

\setmonofont{CourierNew.ttf}[
  Path = ./fonts/,
  UprightFont = CourierNew.ttf
]

% ---------- 额外中文字体族（Fandol） ----------
\setCJKfamilyfont{fs}{FandolFang-Regular.otf}
\setCJKfamilyfont{kai}{FandolKai-Regular.otf}
\setCJKfamilyfont{hei}{FandolHei-Regular.otf}[BoldFont = FandolHei-Bold.otf]
\setCJKfamilyfont{song}{FandolSong-Regular.otf}[BoldFont = FandolSong-Bold.otf]

\newCJKfontfamily\xiaobiaosong{FandolSong-Bold.otf}

% 楷体_GB2312
\newCJKfontfamily\kaifont{FandolKai-Regular.otf}

% 仿宋_GB2312
\newCJKfontfamily\fangsongfont{FandolFang-Regular.otf}

\newCJKfontfamily\kaitifont{FandolKai-Regular.otf}

% 方便你在正文里随时切换：
\newcommand{\fs}{\CJKfamily{fs}}     % 仿宋
\newcommand{\kai}{\CJKfamily{kai}}   % 楷体
\newcommand{\hei}{\CJKfamily{hei}}   % 黑体
\newcommand{\song}{\CJKfamily{song}} % 宋体



